


\begin{frame}
	\frametitle{Réévaluation à base de règles}
	\begin{block}{Idée}
		Dans cette slide, ça doit valoir la peine de repasser un peu de temps sur la
		question des modeleurs paramétriques. %
		Expliquer qu'avec les outils et notre première contribution notre deuxième
		contribution est un modeleur paramétrique à base de règles disposant d'une
		fonctionnalité de rejeu. %
	\end{block}
	À partir de l'enregistrement du processus de construction, nous proposons~:
	\begin{itemize}
		\item un mécanisme de nomination persistante, pour identifier les orbites
		      de manière unique y compris dans le modèle regénéré,
		\item un mécanisme de réévaluation, permettant de regénérer des objets à
		      partir de l'édition d'un processus constructif. %
	\end{itemize}
\end{frame}


% \begin{frame}
% 	\frametitle{Nomination persistante des brins}
% 	\begin{block}{Idée}
% 		Le plan ici, c'est de concrétiser l'approche par historique pour
% 		l'identification des entités topologiques. %
% 		On le fait pour les brins, car c'est générique et sans ambiguïté. %
% 		Par contre, ça devient technique, j'ai besoin de prendre du recul. %
% 	\end{block}
% \end{frame}

% \begin{frame}
% 	\frametitle{Nomination persistante des orbites}
% 	\begin{block}{Idée}
% 		Je ne sais pas trop comment m'y prendre, mais il faudra de toute façon
% 		montrer un DAG (c'est une contribution mise en avant). %
% 		Continuer d'exploiter le cas d'étude et créer le DAG d'évaluation d'une
% 		orbite. %
% 		Montrer qu'à l'aide d'un nom persistant de brin et de l'analyse syntaxique
% 		des opérations, on peut reconstituer l'historique des évolutions d'une
% 		orbite dans un objet. %
% 	\end{block}
% \end{frame}

% \begin{frame}
% 	\frametitle{Réévaluation}
% 	\begin{block}{Idée}
% 		Même problème que la slide précédente, je n'arrive pas à imaginer une
% 		présentation simple. %
% 		Expliquer que l'on peut procéder à la réévaluation, de manière
% 		relativement simple, grâce aux DAG d'évaluation. %
% 		Le principe est simplement de les recréer sous forme de DAG de rejeu qui
% 		prennent en compte l'édition du processus constructif. %
% 	\end{block}
% \end{frame}

\begin{frame}[c,fragile]
	\frametitle{Nommage Persistant de brins}
	\begin{minipage}[t][.4\textheight][t]{\textwidth}
		\begin{overlayarea}{\textwidth}{.25\textheight}
			\centering
			\includegraphics<1-2>[height=.35\textheight]{rule-create-square-face-n3-n6-highlight}%

			\only<1-2>{Création de carré}

			\includegraphics<3-5>[height=.3\textheight]{rule-face-extrusion-compact-n2-n4-highlight}%

			\only<3-5>{Extrusion de face}

			\includegraphics<6-8>[height=.175\textheight]{rule-face-triangulation-3D-n0-highlight}%

			\only<6-8>{Triangulation de face}

			\includegraphics<9-11>[height=.175\textheight]{rule-colour-face}%

			\only<9-11>{Coloration de face}


		\end{overlayarea}
	\end{minipage}
	\begin{columns}[b]
		\begin{column}{.45\textwidth}
			\begin{minipage}[t][.25\textheight][b]{\textwidth}
				\centering
				\includegraphics<2-4>[width=.6\columnwidth]{persistent-naming-step-1}%
				\includegraphics<5>[width=.6\columnwidth]{persistent-naming-step-2}%
				\includegraphics<6-7>[width=.6\columnwidth]{persistent-naming-step-2}%
				\includegraphics<8-10>[width=.6\columnwidth]{persistent-naming-step-3}%
				\includegraphics<11>[width=.6\columnwidth]{persistent-naming-step-4}%
			\end{minipage}
		\end{column}
		\hfill
		\begin{column}{.49\textwidth}
			\begin{minipage}[c][.4\textheight][c]{\textwidth}
				\begin{block}{Spécification paramétrique}
					\begin{overlayarea}{\textwidth}{.2\textheight}
						\ttfamily\footnotesize
						% \begin{semiverbatim}
						\only<1->{1-{\color{mydarkgreenink}square}(pos)}\\
						\only<3->{2-{\color{mydarkgreenink}extrusion}(\alt<3>{6}{[1n6]}, vec)}
						\only<6->{3-{\color{mydarkgreenink}triangulation}(\alt<6>{33}{[1n3;2n2]})}
						\only<9->{4-{\color{mydarkgreenink}colour}(\alt<9>{35}{[1n3;2n4;3n0]})}
						% \end{semiverbatim}
					\end{overlayarea}
				\end{block}
			\end{minipage}
		\end{column}
	\end{columns}
\end{frame}

\begin{frame}
	\frametitle{Nommage Persistant d'une orbite}
	\begin{columns}[b]
		\begin{column}{.45\textwidth}
			\begin{minipage}[t][.25\textheight][t]{\textwidth}
				\begin{overlayarea}{\textwidth}{.25\textheight}
					\centering
					\includegraphics<1>[width=.6\textwidth]{rule-colour-face-n0-highlight}%
					\includegraphics<2>[width=\textwidth]{rule-face-triangulation-3D-n0-highlight-boxed}%
					\includegraphics<3>[width=\textwidth]{rule-face-extrusion-compact-n4-highlight}%
					\includegraphics<4->[width=\textwidth]{rule-create-square-face-n3-highlight}%
				\end{overlayarea}
			\end{minipage}
			\begin{minipage}[t][.5\textheight][b]{\textwidth}
				\begin{figure}
					\includegraphics[width=.6\columnwidth]{persistent-naming-pn3}%
				\end{figure}
			\end{minipage}
		\end{column}
		\hfill
		\begin{column}{.49\textwidth}
			\begin{minipage}[t][.25\textheight][t]{.85\textwidth}
				\ttfamily\scriptsize
				1-{\color{mydarkgreenink}\alt<4->{\color{orange}}{}square}(pos)\\
				2-{\color{mydarkgreenink}\alt<3>{\color{orange}}{}extrusion}([1n6],vec)\\
				3-{\color{mydarkgreenink}\alt<2>{\color{orange}}{}triangulation}([1n3;2n2])\\
				4-{\color{mydarkgreenink}colour}(\alt<1>{{\color{orange}{[1n3;2n4;3n0]}}}{[\alt<4->{{\color{orange}1n3}}{1n3};\alt<3>{{\color{orange}2n4}}{2n4};\alt<2>{{\color{orange}3n0}}{3n0}]})
			\end{minipage}
			\begin{minipage}[t][.5\textheight][b]{\textwidth}
				\begin{figure}
					\includegraphics<1>[height=.5\textheight]{pn7-short-nolevels-step-1}%
					\includegraphics<2>[height=.5\textheight]{pn7-short-nolevels-step-2}%
					\includegraphics<3>[height=.5\textheight]{pn7-short-nolevels-step-3}%
					\includegraphics<4>[height=.5\textheight]{pn7-short-nolevels-step-4}%
					\includegraphics<5>[height=.5\textheight]{pn7-short-nolevels}%
				\end{figure}
			\end{minipage}
		\end{column}
	\end{columns}
\end{frame}

\subsection{Réévaluation}

\begin{frame}
	\frametitle{Réévaluation}
	\begin{columns}
		\begin{column}{.35\textwidth}
			\begin{minipage}[t][.3\textheight][c]{\textwidth}
				\begin{figure}
					\includegraphics<1>[width=.6\columnwidth]{short-1-reevaluation-create-face-darts}%
					\includegraphics<2>[width=.6\columnwidth]{short-1ADD-reevaluation-insert-vertex-darts}%
					\includegraphics<3>[width=.6\columnwidth]{short-2-reevaluation-extrude-face-darts}%
					\includegraphics<4>[width=.6\columnwidth]{short-3-reevaluation-triangulate-face-darts}%
					\includegraphics<5>[width=.6\columnwidth]{short-4-reevaluation-color-face-nodarts}%
				\end{figure}
			\end{minipage}
			%
			\begin{minipage}[b][.7\textheight][t]{\textwidth}
				\begin{figure}
					\includegraphics[height=.5\textheight]{pn7-short-nolevels}%
				\end{figure}
			\end{minipage}
		\end{column}
		\begin{column}{.65\textwidth}
			\hfill
			\begin{minipage}[t][.3\textheight][c]{.65\textwidth}
				\ttfamily\scriptsize
				1-{\color{mydarkgreenink}\alt<1>{\color{orange}}{}square}(pos)\\
				ADD1-{\color{mydarkgreenink}\alt<2>{\color{orange}}{}insert}({\color{red}3})\\
				2-{\color{mydarkgreenink}\alt<3>{\color{orange}}{}extrusion}([1n6], vec)\\
				3-{\color{mydarkgreenink}\alt<4>{\color{orange}}{}triangulation}([1n3;2n2])\\
				4-{\color{mydarkgreenink}\alt<5>{\color{orange}}{}colour}([1n3;2n4;3n0])
			\end{minipage}
			\begin{minipage}[b][.7\textheight][t]{\textwidth}
				\begin{figure}
					\includegraphics<1>[height=.6\textheight]{pn7-reval-short-nolevels-step-1}%
					\includegraphics<2>[height=.6\textheight]{pn7-reval-short-nolevels-step-2}%
					\includegraphics<3>[height=.6\textheight]{pn7-reval-short-nolevels-step-3}%
					\includegraphics<4>[height=.6\textheight]{pn7-reval-short-nolevels-step-4}%
					\includegraphics<5>[height=.6\textheight]{pn7-reval-short-nolevels}%
				\end{figure}
			\end{minipage}
		\end{column}
	\end{columns}
\end{frame}

\begin{frame}
	\frametitle{Configuration de la réévaluation ou stratégies}
	\begin{block}{Idée}
		Je pense que ça peut être intéressant de s'attarder une minute sur les
		possibilités qu'offrent les stratégies pour le rejeu. %
		Ça peut passer par l'illustration d'exemples avec différentes
		configurations comme dans le manuscrit. %
	\end{block}
\end{frame}




%%% Local Variables:
%%% mode: LaTeX
%%% TeX-master: "../../main"
%%% End:
